\begin{titlepage}
  \centering
  \vspace*{2cm}
  {\Large KU Leuven\\Faculty of Engineering Science\\Master of Engineering: Computer Science\par}
  \vspace{2.5cm}
  {\Huge\bfseries Ontwikkeling van een live race-engineeringdashboard voor endurancewedstrijden\par}
  \vspace{2cm}
  {\Large Jeff Mathieu\par}
  \vfill
  \begin{tabular}{ll}
    Opleidingsonderdeel: & Industriële Stage: Computerwetenschappen \\
    Stageplaats: & MotorSport Training Center \\
    Dagelijkse begeleider: & Jan Wouters \\
    Stageperiode: & 22 juni -- 31 juli 2026 (zes weken) \\
    Academiejaar: & 2026--2027 \\
  \end{tabular}
  \vfill
  {\large \today\par}
\end{titlepage}

\pagenumbering{roman}
\chapter*{Abstract}
\addcontentsline{toc}{chapter}{Abstract}
Tijdens deze industriële stage van zes weken werd voor MotorSport Training Center een lokale desktopapplicatie ontwikkeld die live timingdata omzet in bruikbare informatie voor een race-engineer. De bestaande timingwebsites tonen de officiële rangschikking en actuele tijden, maar bieden niet alle teamgerichte vergelijkingen, voorspellingen en strategische hulpmiddelen die tijdens een endurancewedstrijd nodig zijn. De ontwikkelde Electron-applicatie leest daarom periodiek de timingpagina, normaliseert data van GetRaceResults en RIS Timing, bewaart de racehistoriek lokaal en berekent analyses buiten de gebruikersinterface.

Het resultaat is een dashboard voor maximaal drie gelijktijdig gevolgde wagens, met afzonderlijke modi voor race, kwalificatie en vrije training. Het systeem berekent onder andere ronde- en sectorgemiddelden, geldige beste tijden, voorspelde rondetijden, marges tegenover configureerbare referentietijden, vergelijkingen tussen piloten en klassewagens, inhaalindicaties en een pitstopplanning. Neutralisaties, pitgerelateerde rondes en onvolledige data worden expliciet behandeld. Een afzonderlijk grafiekvenster maakt langere trends zichtbaar. De data wordt per sessie opgeslagen in leesbare JSON-, JSONL- en CSV-bestanden, waardoor een sessie na een onderbreking kan worden hervat.

De ontwikkeling verliep iteratief en in nauw overleg met de opdrachtgever. Praktijktests en screenshots brachten verschillende subtiele fouten aan het licht, onder meer bij rondeachterstanden, intervalvelden, pitstopprojecties en wisselende timingproviders. Deze fouten leidden tot modulaire domeinfuncties en een uitgebreide automatische testsuite. Builds voor Windows en macOS worden via GitHub Actions gemaakt en als release gepubliceerd. De definitieve validatie tijdens de geplande test in Spa is op het moment van schrijven nog niet uitgevoerd en wordt daarom in dit verslag duidelijk als toekomstig werk aangeduid.

\textbf{Trefwoorden:} motorsport, live timing, Electron, race-engineering, rondetijdanalyse, pitstopstrategie, softwaretesten.

\chapter*{Administratieve gegevens}
\addcontentsline{toc}{chapter}{Administratieve gegevens}
\begin{tabularx}{\textwidth}{>{\bfseries\raggedright\arraybackslash}p{5.5cm}X}
Naam student & Jeff Mathieu \\
Opleiding & Master in de ingenieurswetenschappen: Computerwetenschappen \\
Stagebedrijf & MotorSport Training Center \\
Dagelijkse stagebegeleider & Jan Wouters \\
Contact begeleider & info@motorsport-trainingcenter.be \\
Stageperiode & 22 juni -- 31 juli 2026, zes weken \\
Opleidingsonderdeel & Industriële Stage: Computerwetenschappen \\
KU Leuven-coördinator & Prof. Tom Van Cutsem \\
\end{tabularx}
